\documentclass[11pt]{article}
\usepackage{fontspec}
\setmainfont{Times New Roman}
\setsansfont{Arial}
\setmonofont{Consolas}
\usepackage{amsmath,amssymb,mathtools}
\usepackage{geometry}
\usepackage{microtype}
\geometry{a4paper,margin=2.05cm}
\setlength{\parindent}{1.1em}
\setlength{\parskip}{0pt}
\newcommand{\foreign}[1]{#1}
\providecommand{\frA}{\mathfrak A}
\providecommand{\frB}{\mathfrak B}
\providecommand{\frC}{\mathfrak C}
\providecommand{\frG}{\mathfrak G}
\providecommand{\frL}{\mathfrak L}
\providecommand{\frM}{\mathfrak M}
\providecommand{\frN}{\mathfrak N}
\providecommand{\frR}{\mathfrak R}
\providecommand{\frS}{\mathfrak S}
\providecommand{\frT}{\mathfrak T}
\providecommand{\frU}{\mathfrak U}
\providecommand{\frako}{\mathfrak o}
\providecommand{\mA}{\mathfrak A}
\providecommand{\mbarA}{\overline{\mathfrak A}}
\providecommand{\Gcal}{\mathcal G}
\providecommand{\Sp}{\operatorname{Sp}}
\providecommand{\glqq}{``}
\providecommand{\grqq}{''}
\emergencystretch=220em
\allowdisplaybreaks
\sloppy
\hbadness=10000
\vbadness=10000
\hfuzz=5pt
\vfuzz=12pt
\raggedbottom
\begin{document}

\section*{40. Nekomutativne algebry}

\begin{center}
\emph{\foreign{Math. Zs. 37 (1933), 514--541}}
\end{center}

\begin{center}
\textbf{Nekomutativna algebra.}\\[0.5em]
Od\\[0.25em]
Emmy Noether v Göttingenu.
\end{center}

Glavne teoremy komutativnoj algebry, kako znano, soderžet se v teoriji Galoisa; jej prědhodi teorija polj odcěpjenja i razpadnyh polj, to jest polj, dostatočnyh za to, aby dano mnogoročje odcěpilo linearny faktor odnosno cělo razpadlo se na linearne faktory.

Sootvětnu čest algebry ja tut razvijam v nekomutativnom, osoblivo v hiperkompleksnom. Pri tom principialno rabotam nekomutativnymi metodami, imenno predstavjenjem v nekomutativnyh tělah. Na koncu pokazu, kako vyšespomenute teoremy komutativnoj algebry možno obosnovati sovsěm paralelno s predstavjenjem v komutativnyh poljah.

Osnovna teorija predstavjenj, kojoj prědhodi kratka teorija avtomorfizmov (§1), jest daljne razvitje toj, ktora opiraje se na moduly predstavjenj.\footnote{E. Noether, \foreign{Hyperkomplexe Größen und Darstellungstheorie}, Math. Zeitschr. 30 (1929), S. 641--692, dalje cit. \foreign{Darstellungstheorie}. Sr. izloženje pri van der Waerden, \foreign{Moderne Algebra II}.} Osoblivo razsmatram pored direktnogo predstavjenja i recipročno predstavjenje, pri čem jedno može se svesti na drugo; ono teper porođaje se posrědstvom recipročnogo modula predstavjenja. Prednost jest v tom, že recipročny modul predstavjenja, to jest dvojny modul oziroma bimodul, možno takodže razuměti kako jednostronny modul nad kolcom razširenja (§2). Tym predstavjenje v nekomutativnyh tělah svodži se na teoriju idealov v kolcu razširenja; nerazložime klasy predstavjenj sootvětsvujut nerazložimym idealnym klasam kolca razširenja, v točnoj analogiji s faktami, ktore pri predstavjenju hiperkompleksnyh sistemov v njihovoj komutativnoj oblasti koeficientov važe za sam sistem (§3 i 4).

Odtud slědujut strukturne teoremy za matrične kolca nad nekomutativnymi těly posrědstvom zamětky (§5), že \emph{vsako podkolco znači predstavjenje tym tělom, odnosno recipročno predstavjenje recipročno izomorfnym tělom}. To recipročno izomorfno tělo jest prvi analog minimalnogo polja odcěpjenja i razpadnogo polja v komutativnom slućaju: ono posrědkuje vse recipročne predstavjenja prvogo stepena i, pri konečnom rangu nad centrom, čto dotolě ne bylo predpostavjeno, takodže polny razpad v direktnu sumu členov ranga \(1\). Iz toga vyhodi Galoisova teorija za těla (§6) i jednočasno vyražaje se fakt (§7), že recipročno izomorfne dělitvene algebry, to jest těla konečnogo ranga nad centrom, porođajut obratne klasy v R. Brauerovoj grupě klasov algeber.

Drugo obosnovanje Galoisovej teorije, obćejše za proste sistemy, jest skoro neposrědnje slědstvo teorema o komutujučih podkolcah, ktory sam skoro direktno nadvezuje se na prědloženno razsmotrenje avtomorfizmov.

Do tega města razsmotrenja idut čisto nekomutativno. Ale i pitanje o \emph{komutativnyh} razpadnyh poljah obrabjaje se nekomutativno posrědstvom predstavjenja razpadnogo polja dělitvenoju algebraju po zamětkě, prědloženoj v §5 (§7). Zaključenje jest vyšespomenuty prěnos na komutativno (§8) i teorija razpadnyh polj za libovoljne sistemy (§9), čem jednočasno daje se svęz s obyčnoju teorijeju predstavjenj v komutativnom.

Na toj teoriji predstavjenj v komutativnom, i na ``iracionalnyh'' faktornyh sistemah, R. Brauer osnoval teoriju razpadnyh polj.\footnote{Sr. našu obču notu: \foreign{Über minimale Zerfällungskörper irreduzibler Darstellungen}. Sitz. Ber. d. Preuß. Ak. d. Wiss. 1927, S. 221--228. Zamětka na S. 222 soderži ``otčet'' o stanju. Glavne teoremy nastali nezavisno i skoro jednočasno. Sr. dalje R. Brauer, \foreign{Über Systeme hyperkomplexer Zahlen}, §3. Math. Zeitschr. 30 (1929), S. 79--107. A. A. Albert pozdnje znovu našel te teoremy nezavisno.} Zaradi togo komutativnogo obosnovanja on, tako kako pozdnje Albert, mora predpostaviti centar kako sovršeno polje, ograničenje nepotrěbno po nekomutativnom obosnovanju. S tym samym ograničenjem, takodže s teorijeju predstavjenj v komutativnom, R. Brauer i K. Shoda, poslě poznanja moje Galoisovej teorije za nekomutativne těla, dalje razvili teoriju: R. Brauer dal vyšespomenuty teorem o komutujučih podkolcah, K. Shoda nezavisno polnu teoriju takodže za poluproste sistemy.\footnote{R. Brauer, \foreign{Über die algebraische Struktur von Schiefkörpern}, §2. Journ. f. Math. 166 (1932), S. 241--252. K. Shoda, \foreign{Über die galoissche Theorie der halbeinfachen hyperkomplexen Systeme}. Math. Ann. 107 (1932), S. 252--258. Te rezultaty razvijal takodže J. Levitzki.}

Nakonec treba spomenuti, že za slućaj těla kratko izloženje nahodi se pri van der Waerden II (§128), v svęzi s mojeju lekcijeju lěta 1928, kogda ja prvi raz govorila o tih vprašanjah. Izloženje van der Waerden prinosa několiko uproščenj proti toj lekciji; jedne ja prevzela v drugoj lekciji (zima 1929/30),\footnote{Prvu lekciju obdelal G. Köthe, drugu M. Deuring.} druge tek tut prevzela i prodolžila. Osoblivo prěnos hiperkompleksnoj dokaznoj metody Galoisovej teorije na komutativno (§8) proishodi tek iz drugoj lekcije, kde nekomutativny dokaz (§6, 3.) byl suščstveno uproščen proti prvoj lekciji. Teorem o komutujučih podkolcah (§5) i na njem osnovana druga dokazna metoda nekomutativnoj Galoisovej teorije (§6, 1. i 2.) dodani sut tek po drugoj lekciji po poznanju R. Brauera (zam. 3) i van der Waerden §128; v §5, 3 sut však suščstveno slabše predpostavjenja konečnosti než tam.

Něktore zaključivanja prvoj lekcije, metoda tvorjenja prěsěkov razlikovyh idealov, sut sice složnějše, ale imajut prednost, že prěnoset se na sistemy beskonečnogo ranga i na cělobrojne sistemy.\footnote{Metoda jest v suščnosti vozproizvedena pri G. Köthe, \foreign{Schiefkörper unendlichen Ranges über dem Zentrum} §5, Math. Ann. 105 (1931), S. 15--39. Sr. takodže zam. 10. Ta metoda prěsěka teper zaměnjena jest skoro trivialnym faktom (§4, 1.), ktory prvi zamětil van der Waerden, že dvostranne idealy naležet k invariantnym modulam. Tym možno vse obosnovati prostějšim razloženjem v direktnu sumu, čto však odpada pri beskonečnom rangu i v cělobrojnom slućaju.} Za komutativne cělobrojne sistemy tako dohodi se do svęzi medžu idealnym diferenciovanjem i diferentoju,\footnote{Sr. predavanje v Pragě, Jahresber. d. Deutsch. Math. Ver. 39 (1930), S. 17 (kosaja paginacija).} k čemu slučajno vrnu se podrobněje.

Glavno priměnjenje teorija razpadnyh polj nahodi v teoriji skrěženih produktov i njihovih faktornyh sistemov, ktore same znovu daju osnovu čislovo-teoretičnyh priměnjenj; tut se o tom dalje ne govori.\footnote{Teorija skrěženih produktov razvijena jest v drugoj lekciji; ona jest, s malymi izměnami, aby tut dano teoriju ne predpostavljati, vozproizvedena pri H. Hasse: \foreign{Theory of cyclic algebras over an algebraic numberfield}, Chap. II. Transactions of the Amer. Math. Soc. 34 (1932), S. 171--214. Izloženje bližje lekciji imaje iziti v otčetu M. Deuringa v \foreign{Ergebnisse der Mathematik}.}

\subsection*{§1. Avtomorfizmy, moduly i dvojne moduly}

Teorija predstavjenj na osnově modulov predstavjenj opiraje se na teoriju kolca avtomorfizmov abelovyh grup s operatorami ili bez njih. Pri tom implicitno osnovane zaključivanja, znači odnosi medžu obrazovanjem i računskymi zakonami, budut za izběganje povtorjenj formulovane v několiko prostyh tvrdženjah.

\paragraph{1. Multiplikativno obrazovanje. Asociativny zakon.}
Nehaj najprvo \(\frG\) bude grupa bez operatorov, \(\frA\) jej absolutna oblast avtomorfizmov, to jest sistem vsih homomorfizmov od \(\frG\) v sebe. \(\frA\) jest multiplikativno zakryta, bo produkt \( \sigma\tau \) definovany jest črez
\[
        g(\sigma\tau)=(g\sigma)\tau,
        \qquad g\in\frG,\quad \sigma,\tau\in\frA                         \tag{1}
\]
i javno zadovolja asociativny zakon.

\(\frG\) postavaje, kako znano, grupa s operatorami, ako dana jest množina \(\frB\) symbolov \(O,H,\ldots\), tak že svęzanja \(gO,gH,\ldots\) daju jednoznačno definirane elementy iz \(\frG\) i porođajut homomorfizmy od \(\frG\) v sebe:
\[
        (gh)O=gO\cdot hO .
\]
Ako oblast operatorov \(\frB\) jest multiplikativno zakryta, togda jednoznačno obrazovanje od \(\frB\) na sootvětnu čest oblasti avtomorfizmov jest multiplikativno-homomorfno togda i samo togda, kogda izpolnjena jest asociativna relacija
\[
        g(OH)=(gO)H, \qquad g\in\frG,\ O,H\in\frB,                       \tag{1a}
\]
Sootvětně definovany jest recipročny homomorfizm, ako ide o lěve operatory.

Bo obrazovanje \(\frB\) na \(\overline{\frB}\) dano jest črez \(g\Theta=g\sigma\) za vse \(g\) iz \(\frG\), togda takodže črez
\[
        (g\Theta)H=(g\sigma)\tau=g(\sigma\tau),
\]
tako že (1a) stavaje nužnym i dostatočnym za multiplikativnu homomorfiju. To, že lěve operatory porođajut recipročny homomorfizm, prihodi odtud, že lěvo zapisane avtomorfizmy \(\sigma^\ast g,\tau^\ast\sigma^\ast g\) treba čitati od desna k lěvu: \(\tau^\ast\sigma^\ast g=g\sigma\tau\), operatory pak vsegda od lěva k desnu.

\paragraph{2. Operatorno-homomorfno obrazovanje.}
Ako \(\frG\) jest grupa s operatorami, operatorna homomorfija definovaje se črez
\[
        (gO)\sigma=(g\sigma)O,                                          \tag{2}
\]
odnosno pri lěvyh operatorah črez
\[
        (Og)\sigma=O(g\sigma).                                          \tag{2*}
\]
to jest črez komutativnu svęzanost odnosno prohodny asociativny zakon. Iz zaključivanja 1., obrazovanja na oblast avtomorfizmov, za dvě razne oblasti operatorov vyhodi:

Ako dane sut dvě oblasti operatorov \(\frB\) i \(\frC\), elementi \(\frC\) porođajut \(\frB\)-avtomorfizmy i jednočasno elementi \(\frB\) porođajut \(\frC\)-avtomorfizmy togda i samo togda, kogda oblasti komutativno svęzane sut s \(\frG\), odnosno kogda izpolnjaje se prohodny asociativny zakon:
\[
        (gO)H=(gH)O,                                                    \tag{2a}
\]
odnosno
\[
        (OH)g=O(Hg).                                                    \tag{2a*}
\]

\paragraph{3. Moduly i dvojne moduly nad kolcami.}
Pravy modul \(\frM\) nad kolcom \(\frR\) jest aditivna abelova grupa s elementami iz \(\frR\) kako pravymi operatorami, pri čem pored asociativnoj relacije (1a) izpolnjajut se i distributivne relacije:
\[
        (g+h)\rho=g\rho+h\rho,\qquad
        g(\rho+\sigma)=g\rho+g\sigma.                                  \tag{3a}
\]
Sootvětny fakt važi za lěve moduly.

Medžu dvojnymi modulami treba razlikovati dva vida. Prave moduly nad dvoma kolcami \(\frR\) i \(\frS\) zovut se dvojne moduly, ako kromě za \(\frR\) i \(\frS\) pojedinačno važečih svęzanj još važi
\[
        (m\rho)\sigma=(m\sigma)\rho .
\]
\(\frR\)-lěve, \(\frS\)-prave moduly zovut se dvojne moduly, ako pridodaje se prohodny asociativny zakon
\[
        \rho(m\sigma)=(\rho m)\sigma .
\]

Značenje tyh svęzanj slěduje iz togo, že v slućaju abelovyh grup absolutna oblast avtomorfizmov tvori kolco, absolutno kolco avtomorfizmov\footnote{Pri neabelovyh grupah ide o \glqq{}obobćeno\grqq{} kolco, sr. rabotu H. Fittinga, ktora ima iziti v Math. Annalen. [Math. Annalen 107, S. 514--542.]}. Podrobno obrazovateljno zaključivanje daje:

Ako oblast operatorov \(\frR\) aditivno zapisanoj abelovoj grupy \(\frM\) jest kolco, togda jednoznačno obrazovanje od \(\frR\) na podmnožinu \(\overline{\frR}\) absolutnogo kolca avtomorfizmov jest kolcovo-homomorfno togda i samo togda, kogda \(\frM\) jest pravy modul nad \(\frR\), to jest kogda (1a) i (3a) izpolnjene sut; pri lěvyh modulah recipročno kolcovo-homomorfno.

Ako \(\frM\) jest pravy modul nad kolcami \(\frR\) i \(\frS\), \(\frM\) stavaje dvojny modul, (2a), togda i samo togda, kogda \(\frR\) možno kolcovo-homomorfno obrazovati na podkolco \(\frS\)-avtomorfnogo kolca i jednočasno \(\frS\) na podkolco \(\frR\)-avtomorfnogo kolca. Ako \(\frM\) jest pravy modul nad \(\frS\), lěvy modul nad \(\frR\), dvojny modul, (2a*), znači homomorfno obrazovanje od \(\frS\) na podkolco \(\frR\)-avtomorfnogo kolca, recipročno homomorfno od \(\frR\) na podkolco \(\frS\)-avtomorfnogo kolca.

Odtud osoblivo slěduje:

1) Ako \(\frR\) jest kolco s jedinicoju, \(\frR\) jest direktno izomorfno svojmu avtomorfnomu kolcu kako \(\frR\)-lěvy modul i recipročno izomorfno svojmu avtomorfnomu kolcu kako \(\frR\)-pravy modul.

2) Ako \(\frR\) jest polno matrično kolco nad vobće nekomutativnym tělom \(A\),
\[
        \frR=\sum c_{ik}A=\sum A c_{ik},
\]
togda \(A\) stavaje recipročno izomorfno avtomorfnomu tělu prostyh pravyh idealov, direktno izomorfno avtomorfnomu tělu prostyh lěvyh idealov.

\paragraph{4. Prěhod od pravogo dvojnogo modula k jednostronnomu modulu nad kolcom produkta.}
Na tom prěhodu suščstveno osnovano jest značenje pravogo dvojnogo modula.

\paragraph{Teorem.}
Ako \(\frM\) jest pravy modul nad kolcom \(\frT\), ktore soderži dva elementno medžusobno komutujuče podkolca \(\frR\) i \(\frS\), togda \(\frM\) možno takodže razuměti kako dvojny modul nad \(\frR\) i \(\frS\). Obratno, ako dan jest dvojny modul nad \(\frR\) i \(\frS\), i obstaja kolco produkta \(\frT\), v ktorom \(\frR\) i \(\frS\) elementno komutujut, togda \(\frM\) možno takodže razuměti kako pravy modul nad \(\frT\), pri čem \(\frT\)-svęzanje prodolžaje dane \(\frR\)- i \(\frS\)-svęzanja.

Prvo tvrdženje slěduje iz toga, že \(\frT\) po definiciji može se homomorfno obrazovati na podkolco \(\overline{\frT}\) absolutnogo kolca avtomorfizmov, pri čem \(\frR,\frS\) prěhodet v elementno komutujuče podkolca \(\overline{\frR},\overline{\frS}\). To ravno znači relaciju (2a), harakterizujuču dvojny modul.

Obratno, ako \(\frM\) dan jest kako \(\frR,\frS\)-dvojny modul (pravy modul), togda \(\frR,\frS\) znovu možno homomorfno obrazovati na elementno komutujuče podkolca \(\overline{\frR},\overline{\frS}\) absolutnogo kolca avtomorfizmov. V absolutnom kolcu avtomorfizmov obstaja za \(\overline{\frR},\overline{\frS}\) kolco produkta\footnote{Kolco produkta znači tut samo kolco, porođeno iz dvuh kolc s obćim nadkolcom; produkt ne mora byti direktny. Že pri danom prěsěku ne vsegda mora obstajati kolco produkta, pokazuje slědujuči priměr: nehaj \(\frako\) bude kolco cělyh čisel, položeno \(\frR=\frako[x]\), \(\frS=\frako[y]\) s \(2x-1=0\), \(2y-1=0\), tako že \(\frako\) jest prěsěk \(\frR\) i \(\frS\), ako medžu \(x\) i \(y\) ne zadano nikako svęzanje. Ale ako \(\frR\) i \(\frS\) ležet v obćem nadkolcu, dohodi: \((2x-1)y-x(2y-1)=0\), togda \(x=y\). Teda ne obstaja kolco produkta s prěsěkom \(\frako\).} \(\overline{\frT}\), ktore zaradi elementnoj komutativnosti sostoji iz vsih elementov formy
\[
        \sum_i \bar r_i\bar s_i+\bar r+\bar s
\]
(ako \(\overline{\frR},\overline{\frS}\) imajut jedinice, dodatne členi \(\bar r,\bar s\) otpadajut). Ako teper obstaja takodže kolco produkta \(\frT\), v ktorom \(\frR\) i \(\frS\) elementno komutujut, ono sootvětně sostoji iz vsih elementov formy \(\sum_i r_i s_i+r+s\). Homomorfizmy
\[
        \frR\to\overline{\frR},\qquad \frS\to\overline{\frS}
\]
možno teda dopolniti črez prirěđenje
\[
        \sum_i r_i s_i+r+s
        \longmapsto
        \sum_i \bar r_i\bar s_i+\bar r+\bar s
\]
do homomorfizma \(\frT\to\overline{\frT}\); teda svęzanje
\[
        m\Bigl(\sum_i r_i s_i+r+s\Bigr)
        =\sum_i (m r_i)s_i+m r+m s
\]
jest jednoznačno i definira po §1, 3. \(\frT\)-modul, ktory obhvača dany \(\frR,\frS\)-modul.

\subsection*{§2. Recipročno i direktno predstavjenje}

\paragraph{1. Moduly linearnyh form.}
Prěhod od teorije razvijene v §1 k teoriji predstavjenj osnovan jest na specializaciji tamtyh modulov na moduly linearnyh form i na razsmotrenju njihovih kolc avtomorfizmov.

\(\frS\)-pravy modul \(\frM\), kde \(\frS\) jest kolco s jedinicoju, zove se modul linearnyh form v \(\frS\), ako \(\frM\) jest direktna suma \(n\) jednočlennyh \(\frS\)-modulov:
\[
        \frM=m_1\frS+\cdots+m_n\frS,
\]
tako že \(m_i\frS\) jest operatorno izomorfen k \(\frS\), znači iz \(m_i s=0\) vsegda slěduje \(s=0\). Iz toga vyhodi, že \(\frS\) možno ne samo homomorfno, ale izomorfno obrazovati na podkolco absolutnogo kolca avtomorfizmov.

\paragraph{Teorem 1.}
Ako \(\frM\) jest modul linearnyh form v \(\frS\), togda \(\frS\)-avtomorfno kolco \(\mA\) od \(\frM\) jest recipročno izomorfno, ne samo homomorfno, kolcu \(\mbarA\) vsih \(n\)-rędnyh matric v \(\frS\).

Libovoljny \(\frS\)-avtomorfizm \(\alpha\) od \(\frM\) polno opisan jest črez prirěđenje
\[
        m_i\longmapsto \overline m_i=m_i\alpha;
\]
bo iz toga po definiciji slěduje
\[
        \sum m_i s_i\longmapsto
        \sum \overline m_i s_i
        =\sum_i (m_i\alpha)s_i
        =\sum (m_i s_i)\alpha.
\]
Ako v matričnom zapisu
\[
        (m_1\alpha,\ldots,m_n\alpha)=(m_1,\ldots,m_n)A,
\]
togda, kogda \(\alpha\) proběgaje vse avtomorfizmy, prirěđenje \(\alpha\mapsto A\) daje jedno-jednoznačnu svęz medžu \(\mA\) i \(\mbarA\). Bo \(\frM\) byl predpostavjen kako modul linearnyh form, \(A\) jednoznačno opreděljena jest črez \(\alpha\), i vsaka matrica \(A\) \(n\)-togo stepena v \(\frS\) porođaje avtomorfizm; dalje slěduje:
\[
        \alpha+\beta\mapsto A+B,\qquad
        \alpha\beta\mapsto BA.
\]
Poslědnje po
\[
        (m_1\alpha\beta,\ldots,m_n\alpha\beta)
        =(m_1,\ldots,m_n)AB
        =(m_1,\ldots,m_n)BA
\]
zaradi \(ms\beta=m\beta s\).

\paragraph{2. Predstavjenje i modul predstavjenja.}
Ako pod recipročnym odnosno direktnym predstavjenjem \(n\)-togo stepena kolca \(\frR\) v \(\frS\) razuměti recipročny odnosno direktny kolcovy homomorfizm od \(\frR\) v podkolco kolca vsih \(n\)-rędnyh matric v \(\frS\), togda teorem 1 možno formulovati takodže:

\paragraph{Teorem 1'.}
Avtomorfno kolco \(\mA\) modula linearnyh form \(n\)-togo stepena v \(\frS\) dopušća izomorfno (věrno) recipročno predstavjenje črez polno matrično kolco \(\mbarA\) vsih \(n\)-rędnyh matric v \(\frS\).

Odtud obične teoremy za direktne predstavjenja dobyvajut takodže formu za recipročne.

\paragraph{Definicija.}
Modul linearnyh form \(\frM\) v \(\frS\) zove se recipročny modul predstavjenja od \(\frR\) v \(\frS\), ako \(\frM\) jest dvojny modul nad \(\frR\) i \(\frS\) i jednočasno \(\frR\)- i \(\frS\)-pravy modul. On zove se direktny modul predstavjenja, ako \(\frM\) jest dvojny modul i jednočasno \(\frR\)-lěvy, \(\frS\)-pravy modul.

\paragraph{Teorem 2.}
Vsaky recipročny odnosno direktny modul predstavjenja porođaje klas ekvivalentnyh recipročnyh odnosno direktnyh predstavjenj od \(\frR\) v \(\frS\), i vse recipročne odnosno direktne predstavjenja porođajut se tak.

Nehaj najprvo \(\frM\) bude recipročny modul predstavjenja. Po §1, 3. togda \(\frR\) možno direktno homomorfno obrazovati na podkolco \(\overline{\frR}\) \(\frS\)-avtomorfnogo kolca od \(\frM\); po §2, Teorem 1' \(\overline{\frR}\) dopušća recipročno (věrno) predstavjenje, ktore togda stavaje recipročnym predstavjenjem od \(\frR\). Ako obratno dano jest recipročno predstavjenje \(\frR\to\frR^*\), ono črez složenje s recipročnym izomorfizmom od \(\frR^*\) k podkolcu \(\overline{\frR}\) \(\frS\)-avtomorfnogo kolca daje direktny homomorfizm od \(\frR\) na \(\overline{\frR}\); teda \(\frM\) stavaje \(\frR\)-modul, imenno dvojny modul, to jest modul predstavjenja po §1, 3. Prěhod k drugim \(\frS\)-bazam daje klas ekvivalentnyh predstavjenj.

Ako ide o direktny modul predstavjenja, \(\frR\) obrazovaje se na \(\overline{\frR}\) recipročno homomorfno. Složenje s recipročnym predstavjenjem od \(\overline{\frR}\) davaje direktno predstavjenje od \(\frR\). Obratno, iz direktnogo predstavjenja slěduje recipročno homomorfno obrazovanje na \(\overline{\frR}\); potom \(\frM\) stavaje direktny modul predstavjenja po §1, 3.

\paragraph{Zamětka.}
Naravno, oba vida predstavjenj možno svesti jedno na drugo: direktno predstavjenje od \(\frR\) jest recipročno predstavjenje recipročnoho kolca k \(\frR\) i obratno. Drugo svodženje sostoji v prěhodu od \(\frS\) k recipročnomu kolcu i zaměně matric transponovanymi. To sootvětsvuje prěhodu od \(\frR\)-lěvogo, \(\frS\)-pravogo modula k \(\frR\)- i \(\frS\)-lěvomu modulu.

\paragraph{3. Prěhod od recipročnogo modula predstavjenja k modulu nad kolcom razširenja.}
Prednost recipročnogo modula predstavjenja opiraje se prěd vsim na možny prěhod k kolcu razširenja po §1, 4. Za tut dany specialny slućaj, \(\frM\) modul linearnyh form po \(\frS\), ide teda o prěhodny teorem za moduly predstavjenj:

\emph{Ako \(\frM\) jest pravy modul nad kolcom \(\frT\), ktore soderži dva elementno komutujuče podkolca \(\frR\) i \(\frS\), tako že \(\frM\) stavaje modul linearnyh form po \(\frS\), togda \(\frM\) možno razuměti takodže kako recipročny modul predstavjenja od \(\frR\) v \(\frS\).}

\emph{Ako obratno dan jest recipročny modul predstavjenja od \(\frR\) v \(\frS\), i obstaja kolco produkta \(\frT\), v ktorom \(\frR\) i \(\frS\) elementno komutujut, togda \(\frM\) možno takodže razuměti kako \(\frT\)-modul tako, že \(\frT\)-svęzanje jest prodolženje danyh \(\frR\)-, \(\frS\)-svęzanj.}

V teoriji predstavjenj suščstveno upotrěbjaje se

\paragraph{Slědstvo iz prěhodnogo teorema za moduly predstavjenj.}
Ako \(\frM\) jest recipročny modul predstavjenja od \(\frR\) v \(\frS\) i jednočasno \(\frT\)-modul s \(\frT\) kako kolcom produkta, togda \(\frR\)-, \(\frS\)-izomorfizm od \(\frM\) k modulu \(\frN\) ravnoznačny jest \(\frT\)-izomorfizmu. \emph{Vsakomu klasu izomorfnyh \(\frT\)-modulov sootvětsvuje teda recipročny klas predstavjenja od \(\frR\) v \(\frS\) i obratno.}

Svęz tyh teoremov s §1, 2. i 3. daje osnovnu za dalje svęz medžu matricami komutujučimi s predstavjenjem i \(\frR\)-, \(\frS\)-avtomorfizmami.

\paragraph{Teorem o komutujučih matricah.}
\emph{Ako \(\frR\to\frR^*\) jest recipročno predstavjenje \(n\)-togo stepena od \(\frR\) v \(\frS\), i \(\frB^*\) znači kolco vsih s \(\frR^*\) elementno komutujučih matric (\(n\)-togo stepena v \(\frS\)), togda \(\frB^*\) daje recipročno-izomorfno predstavjenje \(\frR\)-, \(\frS\)-avtomorfizmov, to jest tyh \(\frS\)-avtomorfizmov, ktore jednočasno sut \(\frR\)-avtomorfizmy, porođajučego modula predstavjenja ili takodže, ako kolco produkta \(\frT\) obstaja, \(\frT\)-avtomorfizmov.}

Nehaj \(\frB\) bude kolco vsih \(\frR\)-, \(\frS\)-avtomorfizmov od \(\frM\). Kako podkolco \(\frS\)-avtomorfnogo kolca \(\mA\), \(\frB\) dopušća recipročno izomorfno predstavjenje v \(\frS\). Po §1, 2. i 3. \(\frB\) sostoji iz vsih takyh avtomorfizmov iz \(\mA\), ktore komutativno svęzane sut s \(\frR\), to jest zadovoljajut relaciju (2a). \(\frB\) sostoji teda iz vsih avtomorfizmov, ktore elementno komutujut s tymi iz \(\overline{\frR}\), kde \(\overline{\frR}\) znači obraz od \(\frR\) v \(\mA\). Poneže pri predstavjenju \(\frB\) i \(\overline{\frR}\) ide o recipročnu izomorfiju, ne samo homomorfiju, to jest dokazovany fakt.

\paragraph{Zamětka.}
\(\frR\)-, \(\frS\)-avtomorfizmy značet prirěđenje \(m_i\mapsto m_i'\) tako, že posrědstvom \(m_i'\) nastavaje to samo predstavjenje \(\frR^*\); pri tom \(m_i'\) ne morajut nužno davati \(\frS\)-bazu od \(\frM\), sootvětně faktu, že matrice iz \(\frB^*\) ne morajut byti obratne. ``Predstavjenje'', ako suma \(\sum m_i'\frS\) ne jest več direktna, treba razuměti v prěnesenom smyslu:
\[
        (m_1',m_2',\ldots,m_n')a=(m_1',\ldots,m_n')A
\]
ostava pravdivo, ale \(A\) ne jest več jednoznačno opreděljeno črez \(a\).

Daljna zamětka k prěhodnomu teoremu jest slědujuča: diagonalne matrice \(E\cdot s\) daju \emph{direktno izomorfno} predstavjenje od \(\frS\), ``identično'' od \(\frS\). Že ide tut o \emph{direktny} izomorfizm, slěduje iz toga, že \(\frS\)-avtomorfno kolco od \(\frM\) soderži podkolco recipročno izomorfno k \(\frS\), čije recipročno predstavjenje tut razsmatra se. Bo vsaky jednočlenny sumand \(m_i\frS\) od \(\frM\) jest k \(\frS\) kako pravomu modulu operatorno izomorfen, jeho \(\frS\)-avtomorfno kolco teda k \(\frS\) recipročno izomorfno (§1, 3., slědstvo 1).

Črez \(\frR\) i \(\frS\) jednoznačno definirano prirěđenje od \(\frT\) k matricam v \(\frS\) jest teda \emph{ne predstavjenje kolca \(\frT\), ale razširenje recipročnogo predstavjenja od \(\frR\) k operatorno-homomorfnomu po \(\frS\):}
\[
        \sum r_i s_i\longmapsto \sum R_i s_i.
\]
Osoblivo recipročno predstavjenje od \(\frR\) jest jednočasno \emph{operatorno-homomorfno vzhodno k v centru \(\frR\) ležečemu prěsěku \([\frR,\frS]\) od \(\frR\) i \(\frS\).}

\subsection*{§3. Moduly v odnosu k tělu}

V slědujučem ide samo o predstavjenja v vobće nekomutativnyh tělah. Zato sobirajut se několiko prostyh teoremov o modulah linearnyh form po tělu.

\paragraph{1. Normalna baza podmodula po danoj bazě polnogo modula.}
Nehaj \(A\) bude tělo i
\[
        \frN=x_1A+\cdots+x_nA
\]
modul linearnyh form ranga \(n\); dalje nehaj
\[
        \frL=z_1A+\cdots+z_\ell A
\]
podmodul ranga \(\ell\le n\). \(z_i\) zovut se normalna baza po \(x\), ako sut formy
\[
        z_i=x_i-(x_{\ell+1}\alpha_{i,\ell+1}+\cdots+x_n\alpha_{i,n}),
        \qquad i=1,\ldots,\ell.
\]
Vsaky modul \(\frL\) ima po podobnom numerovanju \(x_i\) normalnu bazu. Bo po podobnom numerovanju dohodi
\[
        \frN=\frL+x_{\ell+1}A+\cdots+x_nA,
\]
teda
\[
        x_i\equiv x_{\ell+1}\alpha_{i,\ell+1}+\cdots+x_n\alpha_{i,n}
        \pmod{\frL},\qquad i=1,\ldots,\ell,
\]
i sootvětny
\[
        z_i=x_i-(x_{\ell+1}\alpha_{i,\ell+1}+\cdots+x_n\alpha_{i,n})
\]
ležet v \(\frL\) i izčerpajut \(\frL\).

\paragraph{2. Modul razširenja.}
Nehaj znovu \(A\) bude tělo, \(P\) podtělo. \(A\)-modul \(N\) ranga \(n\) zove se modul razširenja \(P\)-modula \(M\); \(N=M_A\), ako \(M\) jest podmodul jednakogo ranga v \(N\). Ako
\[
        M=y_1P+\cdots+y_nP,
\]
teda
\[
        N=M_A=y_1A+\cdots+y_nA.
\]
Obratno k vsakomu \(P\)-modulu \(M=y_1P+\cdots+y_nP\) obstaja do modulnoj izomorfije jednoznačno modul razširenja \(M_A\). Bo modul \(y_1A+\cdots+y_nA\) soderži k \(M\) izomorfny podmodul \(M\), ktory može byti identificiran s \(M\).

\paragraph{Slědstvo a).}
Ako \(z_1,\ldots,z_s\) sut vzhodno k \(P\) linearno nezavisne elementy iz \(M\), togda oni sut takodže linearno nezavisne po \(A\) v \(N=M_A\), bo jih možno dopolniti do bazy od \(M\) i tym samym od \(N\). Vsaky \(P\)-modul
\[
        T=z_1P+\cdots+z_sP
\]
iz \(M\) porođaje teda modul razširenja
\[
        T_A=z_1A+\cdots+z_sA\subset M_A.
\]

\paragraph{Slědstvo b).}
Vsaky \(P\)-modul \(T=z_1P+\cdots+z_sP\) iz \(M\) jest suženny modul, to jest prěsěk svojego modula razširenja \(T_A\) s \(M\):
\[
        T=T_A\cap M .
\]
Bo \(T\subseteq T_A\cap M\); leži li obratno \(a\) v \(T_A\cap M\), dohodi
\[
        a=\sum z_i\alpha_i,
\]
teda elementi \(M\): \(a,z_1,\ldots,z_s\) sut zavisne po \(A\), teda po slědstvu a) takodže po \(P\); zato \(a\) leži v \(T\).

\paragraph{3. Teorem o invariantnyh modulah.}
Nehaj znovu \(N=M_A\) jest modul razširenja \(P\)-modula \(M\) ranga \(n\), dalje nehaj \(P\) bude polno invariantno tělo vzhodno k grupě \(\frG\) kolcovyh avtomorfizmov od \(A\). Grupa \(\frG\) definovaje se kako oblast operatorov od \(N=M_A\) črez postavjenja
\[
        Gm=m\quad\text{za }m\text{ iz }M\text{ i }G\text{ iz }\frG
\]
i
\[
        G\sum m_i\alpha_i=\sum m_i\,G(\alpha_i)
        \quad\text{za }m_i\text{ iz }M\text{ i }\alpha_i\text{ iz }A.
\]

\paragraph{Lema.}
Po tyh postavjenjah \(M\) sostoji iz vsih elementov od \(N\), ktore ostavajut nepodvižne pri \(\frG\). Bo ako \(x_1,\ldots,x_n\) jest \(P\)-baza od \(M\), teda \(w=x_1\alpha_1+\cdots+x_n\alpha_n\), iz \(G(w)=w\) za vsako \(G\) iz \(\frG\) slěduje, že takodže \(G(\alpha_i)=\alpha_i\), i po predpostavjenju \(\alpha_i\) leži v \(P\).

\paragraph{Teorem.}
Moduly razširenja \(L_A\) podmodulov \(L\) od \(M\), i samo oni, sut dopustime podmoduly vzhodno k \(\frG\) kako oblasti operatorov.

Prva polovica jest jasna. Obratno, nehaj \(L\) bude dopustim vzhodno k \(\frG\), nehaj \(z_1,\ldots,z_\ell\) bude normalna baza \(L\) (sr. 1.) po \(x_1,\ldots,x_n\), gde \(x_1,\ldots,x_n\) izbrana sut kako \(P\)-baza od \(M\), teda \(\frG\) elementno dopušća jih. Po predpostavjenju
\[
        G(z_i)=x_i-\bigl(x_{\ell+1}G(\alpha_{i,\ell+1})
              +\cdots+x_nG(\alpha_{i,n})\bigr)
\]
jest element od \(L\) za vsako \(G\) iz \(\frG\), teda pri fiksnom \(G\) linearno vyražimy črez \(z\); iz sravnjenja koeficientov pri \(x_1,\ldots,x_n\) dostava se \(G(z_i)=z_i\) za vsako \(G\) iz \(\frG\). Teda po lemi \(z\) ležet v \(M\), i \(L\) stavaje razširenje \(P\)-modula
\[
        T=z_1P+\cdots+z_\ell P
\]
iz \(M\), pri čem \(T=L\cap M\) (po 2).\footnote{Teoremy togo paragrafa ostavajut po prostyh dobro-uporędkovyh zaključivanjah validne i ako rang od \(M\) po \(P\) stane beskonečny (sr. G. Köthe: \foreign{Ein Beitrag zur Theorie der kommutativen Ringe ohne Endlichkeitsvoraussetzung} §1, Gött. Nachr. 1931, S. 195--207).}

\subsection*{§4. Hiperkompleksne sistemy i njihove klasy predstavjenj}

\paragraph{1. Teorem razširenja.}
Teper svęzujemo teoremy predstavjenj iz §2 s modulnymi teoremami iz §3. Na město obćih kolc \(\frR\) razsmatramo hiperkompleksne sistemy s jedinicoju, \(S,T,\ldots\), nad komutativnym poljem \(P\), to jest
\[
        S=x_1P+\cdots+x_nP;
\]
na město libovoljnyh kolc predstavjenja \(\frS\) samo těla \(A,B,\ldots\), i to, ako ništo drugo ne zaměčeno, taka s centrom \(P\). V tom slućaju vsegda obstaja kolco produkta, v ktorom \(S\) i \(A\) elementno komutujut, s prěsěkom \(P\); ide o direktny produkt \(S\times A\), ktory jednočasno stavaje modul razširenja \(S_A\) od \(S\) v smyslu §3,
\[
        S_A=x_1A+\cdots+x_nA,
\]
i zato zove se takodže kolco razširenja.

Pri toj specializaciji teorem o invariantnyh modulah (§3, 3.) prěhodi v

\paragraph{Teorem razširenja.}
Vsaky dvostronny ideal iz \(S_A\) jest ideal razširenja ideala iz \(S\), imenno razširenje svojego prěsěka s \(S\).

Bo ako osnovna grupa \(\frG\) jest grupa vnutrišnjih avtomorfizmov od \(A\), \(P\) kako centar od \(A\) stavaje polno invariantno tělo, medžutym elementna komutativnost \(A\) s \(S\) znači, že \(\frG\) po §3, 3. razširjena jest na \(S_A\) tako, že \(S\) stavaje polna oblast invariantov. Vsaky dvostronny ideal \(\mathfrak a\) iz \(S_A\) jest dopustima podgrupa, zaradi \(\alpha^{-1}\mathfrak a\alpha\subseteq\mathfrak a\), teda po §3, 3. razširenje svojego prěsěčnogo modula \(\mathfrak a\cap S\), ktory stavaje ideal v \(S\).

Dodatok: Ako \(S\) jest komutativno, \(S\) stavaje centrom od \(S_A\). Bo \(S\) leži v centru i jest polna invariantna oblast vzhodno k vnutrišnjim avtomorfizmam porođenym od \(A\). Obćejše, centar od \(S\) stavaje takodže centrom od \(S_A\).

Označenje: Pod \(A_r,B_n,\ldots\) treba vsegda razuměti matrične kolca naznačenogo stepena nad \(A,B,\ldots\), teda
\[
        A_r=\sum A c_{ik}=\sum c_{ik}A.
\]

V slědujučem suščstveno upotrěbjaje se slědstvo, teorem razširenja za proste sistemy: Ako \(S\) jest prosty sistem\footnote{Prosty sistem znači kako obyčno \glqq{}dvostronno prosty\grqq{}.}, togda \(S_A\) takodže stavaje prosty:
\[
        S_A=\sum B c_{ik}=\sum c_{ik}B=B_t
\]
s prirěđenym tělom \(B\), čiji centar obhvača \(P\). Pri tom \(A\), i tym takodže \(B\), može byti beskonečnogo ranga nad \(P\); samo \(S\) predpostavja se kako hiperkompleksny. Ako \(S\) i \(A\) imajut obći centar \(P\), togda \(P\) stavaje takodže centrom od \(S_A\).

Obćejše važi još: ako \(S\) jest prosty hiperkompleksny sistem, togda takodže direktny produkt \(S\times A_r\) jest prosty. Bo \(S\times A_r\) sovpada s \(S_A\times P_r\), teda stavaje \(B_{tr}\), ako \(S_A=B_t\).

\paragraph{2. Klasy predstavjenj.}
Pod recipročnym ili direktnym predstavjenjem hiperkompleksnogo sistema treba, kako obyčno, razuměti tako, ktore jest operatorno-homomorfno vzhodno k oblasti koeficientov \(P\) (§2, konec) i različno od nulovogo predstavjenja. Svęz prěhodnogo teorema i jego slědstva (§2, 3.) s teoremom razširenja iz 1. daje

\paragraph{Teorem o klasah predstavjenj.}
Ako \(S\) jest hiperkompleksny s jedinicoju, s \(P\) kako oblastju koeficientov, a \(A\) tělo s centrom \(P\), togda \(S\) ima toliko raznyh nerazložimyh recipročnyh klasov predstavjenj od \(S\) v \(A\), koliko ima klasov prostyh pravyh idealov v ostatkovom kolcu od \(S_A\) po njegovom radikalu. Ako \(S\) jest prosty sistem, togda ima točno jeden nerazložimy recipročny i jeden nerazložimy direktny klas predstavjenja v \(A\).\footnote{Sr. \foreign{Darstellungstheorie}, zam. 20), za slućaj, že \(A\) jest prirěđeno avtomorfno tělo od \(S\).}

Bo po slědstvu iz prěhodnogo teorema (§2, 3.) proste recipročne klasy predstavjenj i proste modulne klasy po \(S_A\) sootvětsvujut jedno-jednoznačno. Poneže \(S_A\) jest konečnogo ranga po tělu \(A\), ono zadovolja maksimalno i minimalno uslovje za jednostronne idealy, teda prvi děl tvrdženja slěduje direktno iz \foreign{Darstellungstheorie}, §19. Ako \(S\) jest prosty, togda takodže \(S_A\) po teoremu razširenja. Zaradi maksimalnogo i minimalnogo uslovja ono jest jednočasno polno reducibilno, ima teda samo jeden prosty jednostronny idealny klas; to jest, \(S\) ima samo jeden nerazložimy recipročny klas predstavjenja v \(A\). S \(S\) jest takodže recipročno izomorfny sistem prosty, ima samo jeden nerazložimy recipročny klas predstavjenja; zato \(S\) ima takodže samo jeden prosty direktny klas predstavjenja.

\paragraph{3. Rangova relacija.}
Fakt, že ako \(S\) jest prosty,
\[
        S_A=\sum Bc_{ik}
\]
jest konečnogo ranga kako po \(A\), tako po \(B\), daje rangovu relaciju:
\[
\begin{gathered}
        n=rt,\qquad n=(S:P),\qquad t^2=(S_A:B),\\
        r=\text{stepen nerazložimogo recipročnogo predstavjenja}.
\end{gathered}                                                            \tag{1}
\]

Bo \(S_A\) samo jest recipročny modul predstavjenja za \(S\) v \(A\) (predstavjenje leži už v \(P\)). Kako matrično kolco nad \(B\), \(S_A\) razpadaje se na \(t\) prostyh pravyh idealov, teda \(S_A\)-modulov, čiji rang po \(A\) jest stepen \(r\) nerazložimogo recipročnogo predstavjenja. Poneže \(n\) jest jednočasno rang od \(S_A\) po \(A\), slěduje rangova relacija.

\paragraph{4. Poostrenje rangovoj relacije, ako \(A\) jest konečnogo ranga po \(P\).}
V tom slućaju važe tri relacije:
\[
        (S:P)=n=rt,                                                     \tag{1}
\]
\[
        (B:P)\,t=(A:P)\,r,                                               \tag{2}
\]
\[
        (S:P)(B:P)=(A:P)\,r^2.                                           \tag{3}
\]
Pri tom (2) znači rang po \(P\) prostogo pravogo ideala iz \(S_A\), vyraženy po 3. črez rang po \(B\) odnosno \(A\), medžutym (3) znači rang prostogo pravogo ideala v matričnom kolcu \(B_n\) i slěduje iz (2) črez množenje s \(r\) po (1).

\subsection*{§5. Priměnjenje teoremov predstavjenja na matrične kolca}

V §4 razsmatrjano recipročno predstavjenje od \(S\) v \(A\) možno takodže razuměti kako direktny homomorfizm od \(S\) v podkolco od \(A_r\), gde \(\overline A\) znači k \(A\) recipročno izomorfno tělo. Princip slědujučih paragrafov jest obratno svesti izslědovanje matričnyh kolc \(A_r\) na teoriju predstavjenj v recipročno izomorfnom tělu \(A\).

\paragraph{1. Reducibilno i nerazložimo vloženje v \(A_r\).}
Nehaj \(A\) i \(\overline A\) sut recipročno-izomorfne těla konečnogo ili beskonečnogo ranga nad svojim centrom \(P\).\footnote{Obće recipročne těla ili sistemy budut označene sootvětnymi grčskymi i latinskymi bukvami.} Prosty sistem \(S\) nad \(P\) zove se nerazložimo odnosno reducibilno vložimy v \(A_r\), ako \(S\) ima nerazložimo odnosno reducibilno recipročno predstavjenje \(r\)-togo stepena v \(A\). Ako \(S\) nerazložimo vložimy jest v \(A_r\), togda on reducibilno vložimy jest v vsih matričnyh kolcah \(A_{rs}\) s \(s>1\), bo ima reducibilne predstavjenja tih i samo tih stepenov (§4, 2.). \(S\) stavaje teda direktno izomorfny někojim podkolcam od \(A_r\) i \(A_{rs}\). Izomorfizm jest prodolženje identičnogo od \(P\);\footnote{Kako už pri predstavjenjah, pri v slědujučem nastavajučih izomorfizmah vsegda ide o take, ktore sut prodolženje identičnogo od \(P\), izomorfne po \(P\), i bez osobnogo spomenanja.} po §4, 3. pri tom \(r\) jest dělitelj ranga \(n\) od \(S\) po \(P\).

Na toj način dobyvajut se vse \(P\) obhvačajuče proste podkolca konečnogo ranga po \(P\) raznyh \(A_f\) (\(f=1,2,\ldots\)); bo vsako tako podkolco jest recipročno izomorfno podkolcu od \(\overline A_f\), dopušča teda reducibilno ili nerazložimo recipročno predstavjenje v \(A\).

\paragraph{2. Teorem o vnutrišnjih avtomorfizmah.}
Ako \(S^{(1)}\) i \(S^{(2)}\) sut dva \(P\) obhvačajuče proste podkolca od \(A_f\), konečnogo ranga po \(P\), i \(S^{(1)}\) i \(S^{(2)}\) izomorfne sut po \(P\), togda toj izomorfizm porođaje se vnutrišnjim avtomorfizmom od \(A_f\).

Bo \(S^{(1)}\) i \(S^{(2)}\) značet, vobće reducibilne, vloženja prostogo sistema \(S\) v \(A_f\), teda direktne predstavjenja \(f\)-togo stepena v \(A\); oni naležet zato k tomu samomu reducibilnomu direktnomu klasu predstavjenj v \(A\); prěnosna matrica porođaje avtomorfizm.

Ako samo \(A\) jest konečnogo ranga po \(P\), teda hiperkompleksno, teorem prěhodi v znany: dva proste podsistemy od \(A_f\), obhvačajuče centar \(P\), izomorfne po \(P\), prěhodet jedna v drugu črez vnutrišnji avtomorfizm od \(A_f\). Osoblivo vsaky avtomorfizm od \(A_f\) jest vnutrišnji.\footnote{To ne važi več, ako \(A\) jest beskonečnogo ranga po \(P\) (sr. Köthe v rabotě citovanoj v zam. 5).}

\paragraph{3. Teorem o elementno komutujučih podkolcah.}
Toj teorem slěduje iz teorema o komutujučih matricah iz §2, 3., znovu po principu spomenutom na počętku paragrafa:

Ako \(S\) jest prosty, \(P\) obhvačajuči sistem iz \(A_f\), i \(R\) znači vsě elementy iz \(A_f\), ktore elementno komutujut s \(S\), togda takodže \(R\) jest matrično kolco: \(R=\overline B_s\). Prirěđene těla \(B\) od \(S_A\) i \(\overline B\) od \(R\) sut recipročno izomorfne. Prěsěk od \(R\) i \(S\) stavaje centrom od \(S\). \(R\) jest tělo togda i samo togda, kogda \(S\) jest nerazložimo vloženo v \(A_f\).

Bo po §2, 3. \(R\) jest direktno izomorfno avtomorfnomu kolcu recipročnogo modula predstavjenja \(\frM\) od \(S\) v \(A\); to pak jest avtomorfno kolco od \(\frM\), razuměnogo kako \(S_A\)-modul. Ako \(\frM\) razpadaje se na \(s\) prostyh \(S_A\)-modulov, i, kako v §4, 1., položeno jest \(S_A=\sum Bc_{ik}\), togda \(R\) jest izomorfno k \(\sum \overline Bc_{ik}\), gde \(B\) i \(\overline B\) sut recipročno izomorfne (bo \(B\) jest recipročno izomorfno avtomorfnomu tělu prostyh pravyh idealov iz \(S_A\), teda prostyh \(S_A\)-modulov po koncu §1, 3.). \(R\) stavaje tělo, \(R=\overline B\), togda i samo togda, kogda \(s=1\), teda \(S\) nerazložimo vloženo jest. Že prěsěk od \(R\) i \(S\) jest centar od \(S\), slěduje neposrědnje iz definicije od \(R\).

\paragraph{4. Poostreny komutantny teorem pri hiperkompleksnom matričnom kolcu.}
V tom slućaju rangove relacije iz §4, 4. daju poostrenje:

Proste, \(P\) obhvačajuče podsistemy iz \(A_f\), gde \(A\) jest konečnogo ranga po svojem centru \(P\), razpadajut se na pary \(S,\overline S\) tako, že \(\overline S\) sostoji iz vsih s \(S\) elementno komutujučih elementov i obratno. Prirěđene těla od \(S_A\) i \(\overline S\) sut recipročno izomorfne, takodže te od \(S\) i \(\overline S_A\). Prěsěk od \(S\) i \(\overline S\) ravny jest obćemu centru. Jedno podkolco jest tělo togda i samo togda, kogda drugo nerazložimo vloženo jest. Produkt rangovyh čisel po \(P\) od \(S\) i \(\overline S\) ravny jest rangu od \(A_f\). Ako \(f=rs=\overline r\,\overline s\), gde \(S\) nerazložimo vložimo jest v \(A_r\) i \(\overline S\) nerazložimo vložimo v \(A_{\overline r}\), prirěđene těla od \(S\) i \(\overline S\) stavajut izomorfne paru komutujučih podkolc iz \(A_g\) s \(f=g s\overline s\).

Suščstveno treba dokazati samo tvrdženje o produktu rangovyh čisel. Nehaj najprvo \(S\) nerazložimo vloženo, teda \(\overline S\) tělo i recipročno izomorfno k \(B\), gde \(S_A=B_t\). Rangova relacija (3) iz §4, 4. daje produktno tvrdženje, bo teper \(f=r\), i rang od \(\overline S\) ravny jest rangu od \(B\). Ako \(S\) reducibilno vloženo, teda \(\overline S=B_s\), togda \(f=rs\), množenje (3) s \(s^2\) daje tvrdženje. Odtud slěduje neposrědnje, že \(\overline S\) sostoji iz vsih elementov elementno komutujučih s \(S\), i tym vse dalje iz 3. kromě poslědnjogo tvrdženja. To vyhodi črez dvojny prěhod k nerazložimomu vloženju. V \(A_r\), s \(f=rs\), \(S\) nerazložimo vloženo jest, \(S\) i \(R\) stavajut parno komutujuče sistemy, pri čem \(R\) izomorfno tělu \(B\). Dalje \(S\), zaradi recipročnosti od \(S\) i \(\overline S\), ravno jest \(\overline B_s\), gde \(\overline B\) ima to samo značenje za \(\overline S\), ktore \(B\) ima za \(S\); teda \(R\) reducibilno vloženo jest v \(A_r\), nerazložimo v \(A_g\) s \(r=g\overline s\). Tut parno komutujuče sistemy stavajut izomorfne k \(B\) i \(\overline B\).

\subsection*{§6. Galoisova teorija prostyh sistemov}

Poostreny komutantny teorem iz §5, 4. vyražaje Galoisovu teoriju prostyh sistemov vzhodno k centru kako osnovnoj oblasti. Razsmatramo najprvo těla, potom obće proste sistemy.

\paragraph{1. Galoisova teorija těla konečnogo ranga po centru.}
Galoisova grupa \(\Gcal\) od \(A\) definovana jest kako grupa vsih avtomorfizmov, ktore sut prodolženje identičnogo od centra \(P\), teda po §5, 2. kako grupa vsih vnutrišnjih avtomorfizmov. Zato
\[
        \Gcal \simeq A^*/P^*,
\]
gde \(A^*\) odnosno \(P^*\) značet multiplikativnu grupu nenulovyh elementov iz \(A\) odnosno \(P\). Podgrupam \(\mathfrak H\) od \(\Gcal\) sootvětsvujut teda jedno-jednoznačno \(P^*\) obhvačajuče podgrupy \(H^*\) od \(A^*\), po
\[
        \mathfrak H \sim H^*/P^* .
\]
Podgrupa \(\mathfrak H\) od \(\Gcal\) zove se zakryta, ako \(H^*\) po pridavanju nuly prěhodi v tělo \(H\). Fakt, že \(C\) dopušća grupu \(\mathfrak H\), ravnoznačny jest tomu, že \(C\) elementno komutuje s \(H\). Tym komutantny teorem §5, 4. prěhodi v

\paragraph{Glavny teorem Galoisovej teorije nekomutativnyh těl.}
Těla \(C\) medžu \(P\) i \(A\) i zakryte podgrupy \(\mathfrak H\) od \(\Gcal\) možut se jedno-jednoznačno prirěditi tako, že \(C\) stavaje polno invariantno tělo od \(\mathfrak H\) i \(\mathfrak H\) polna invariantna grupa od \(C\).

\paragraph{2. Galoisova teorija prostyh sistemov.}
Ako \(A_f\) jest prosty sistem s centrom \(P\), Galoisova grupa \(\Gcal\) znovu jest grupa vnutrišnjih avtomorfizmov, teda izomorfna k
\[
        A_f^*/P^*,
\]
gde teper \(A_f^*\) znači multiplikativnu grupu regularnyh elementov iz \(A_f\). Podgrupa \(\mathfrak H\sim H^*/P^*\) od \(\Gcal\) zove se prosto-zakryta, ako iz \(H^*\) porođeno podkolco \(H\) jest prosty sistem i ako \(H^*\) sostoji iz vsih regularnyh elementov od \(H\). Prěhod od komutantnogo teorema k Galoisovej teoriji tut dan jest črez

\paragraph{Lema.}
Vsaky \(P\) obhvačajuči prosty sistem iz \(A_f\) porođaje se grupu \(H^*\) svojih regularnyh elementov. To jest jasno, ako \(P\) ima beskonečno mnogo elementov; bo s neizvěstnymi stvoreny \glqq{}obći element\grqq{} jest regularny, i črez podobnu specializaciju možno tvoriti regularne bazne elementy po \(P\). Lema važi po Shodi obće.\footnote{Shoda, rabota citovana v zam. 3). Tam izběgaje se vvedenje neizvěstnyh.}

Po lemi komutativnost od \(S\) s prostym sistemom \(\mathfrak S\) znovu ravnoznačna jest tomu, že \(S\) dopušća avtomorfizmy inducirane regularnymi elementami od \(\mathfrak S\). Teda i tut komutantny teorem §5, 4. prěhodi v

\paragraph{Glavny teorem Galoisovej teorije prostyh sistemov.}
Proste, \(P\) obhvačajuče sistemy \(S\) iz \(A_f\) i prosto-zakryte podgrupy \(\mathfrak H\) od \(\Gcal\) možut se jedno-jednoznačno prirěditi tako, že \(S\) stavaje polna invariantna oblast od \(\mathfrak H\) i \(\mathfrak H\) polna invariantna grupa od \(S\).

\paragraph{3. Dokaz posrědstvom principa prodolženja.}
Za slućaj těla dajem još drugi dokaz glavnogo teorema, osnovany na principu prodolženja izomorfizmov, to jest predstavjenj prvogo stepena, i, bo ne upotrěbjaje rangovogo sčitanja, potrebuje menje predpostavjenj konečnosti. Osoblivo tako vidi se směr, v ktorom bude možny prěnos na těla \(S\) beskonečnogo ranga. Dokaz takodže ne upotrěbjaje fakt, že obstaja samo jeden nerazložimy recipročny klas predstavjenja; zato važi točno tako v komutativnom, sr. §8, 2. Prědlagam několiko lem.

\paragraph{Predpostavjenja.}
Prosty sistem \(S\) nad \(P\) nehaj dopušća recipročno predstavjenje prvogo stepena v \(A\), teda jest tělo; pri tom \(A\) može byti konečnogo ili beskonečnogo ranga nad svojim centrom \(P\). Razloženje od \(S_A\) v \(n\) prostyh, operatorno izomorfnyh pravyh idealov po §4, 3. dano nehaj bude črez
\[
        S_A=r_1+\cdots+r_n
            =e_1S_A+\cdots+e_nS_A
            =e_1A+\cdots+e_nA ;
\]
črez \(e_i\) porođeno recipročno predstavjenje definirano jest črez
\[
        e_i\sigma=e_i\sigma_i .
\]

\paragraph{Lema 1.}
\(n\) črez \(e_1,\ldots,e_n\) porođenyh recipročnyh predstavjenj sut razne, teda različne predstavjenja togo samogo klasa predstavjenj. \(S\) ima teda najmenej toliko različnyh predstavjenj, koliko jest jego rang.

Za dokaz razširjamo predstavjenja iz \(S\) po koncu §2 na prirěđenja od \(S_A\), ktore stavajut operatorno-homomorfne po \(A\). Ta prirěđenja definirana sut tut črez
\[
        e_i w=e_i\omega
\]
s \(\omega\in A\) za vsako \(w\in S_A\). Ako by \(e_i\) i \(e_j\), \(i\ne j\), porođali to samo predstavjenje, dohodi by \(e_iw=e_jw\) za vsako \(w\in S_A\). To nemožno zaradi \(e_ie_i=e_i\) i \(e_je_i=0\).

\paragraph{Lema 2.}
Ako \(T\) jest tělo medžu \(P\) i \(S\), \(s\) rang od \(S\) po \(T\), togda vsaky recipročny izomorfizm od \(T\) v \(A\) dopušča najmenej \(s\) različnyh prodolženj.

Recipročny izomorfizm, to jest recipročno predstavjenje od \(T\) v \(A\), nehaj bude posrědkovany razloženjem
\[
        T_A=E_1T_A+\cdots+E_hT_A
            =E_1A+\cdots+E_hA,
        \qquad E_it=E_i\tau ,
\]
gde \(E_i\) sut idempotenty. To ne jest ograničenje obćnosti, bo predstavjenje porođeno baznym elementom \(E_i\alpha\) porođaje se takodže črez \(\alpha^{-1}E_i\alpha\), teda črez idempotent. Pravo množenje s \(S\) daje
\[
        S_A=E_1S_A+\cdots+E_hS_A .
\]
Pri tom vsě \(E_iS_A\) sut ranga \(s\) po \(A\). Bo rang od \(E_iS_A\) jest najvyše \(s\), kako pokazuje vloženje \(T\)-bazy v \(S\) posrědstvom komutativnosti \(S\) i \(A\):
\[
        S=Tw_1+\cdots+Tw_s,
\]
\[
        E_iS_A=E_iT_Aw_1+\cdots+E_iT_Aw_s
              =E_iw_1A+\cdots+E_iw_sA .
\]
Poneže suma \(S_A\) jest ranga \(n=hs\), vsako \(E_iS_A\) ima točno rang \(s\). Teda
\[
        E_iS_A=r_{i1}+\cdots+r_{is}
              =e_{i1}A+\cdots+e_{is}A,
\]
gde \(s\) predstavjenj porođenyh črez \(e_{i1},\ldots,e_{is}\) po lemi 1 vsě stavajut razne. Vsě sut prodolženja predstavjenja od \(T\) porođenogo črez \(E_i\); bo iz \(E_it=E_i\tau\) slěduje
\[
        (e_{i1}+\cdots+e_{is})t=(e_{i1}+\cdots+e_{is})\tau
\]
i tym \(e_{ij}t=e_{ij}\tau\) zaradi direktnogo sumarnogo razloženja.

\paragraph{Definicija.}
Črez \(T\) inducirano razděljenje na klasy \(\mathfrak I_T\) recipročnyh izomorfizmov \(\mathfrak I\) od \(S\) na \(A\) definovano jest tako: vse i samo te izomorfizmy iz \(\mathfrak I\) razsmatrajut se kako ekvivalentne, ktore na \(T\) inducirajut toj samy izomorfizm.

\paragraph{Glavny teorem.}
Ako \(\mathfrak I_T\) jest črez \(T\) inducirano razděljenje na klasy, togda \(T\) jest maksimalno podtělo vzhodno k tomu razděljenju.

Bo ako \(Z\) jest medžutělo od \(T\) i \(S\), stepena \(l\) po \(T\), togda po lemi 2 vsaky klas iz \(\mathfrak I_T\) razpadaje se najmenej na \(l\) klasov iz \(\mathfrak I_Z\). Prěhod od toj formy glavnogo teorema k obyčnoj nastavaje tak, že množina izomorfizmov \(\mathfrak I\) črez množenje svojih elementov s obratnym někojego fiksnogo elementa prěhodi v grupu avtomorfizmov, pri čem klas iz \(\mathfrak I_T\) prěhodi v invariantnu grupu od \(T\). Tvrdženje, že zakryte podgrupy sut polne invariantne grupy, tut je vključeno; treba samo za \(\mathfrak H\sim H^*/P^*\) vzęti tělo \(H\) kako medžutělo \(T\).

\subsection*{§7. Klas podobnyh algeber. Razpadne polja}

Odtud dalje ide samo o tělah \(A,\overline A\) i matričnyh kolcah \(A_f\), ktore sut konečnogo ranga po svojem centru \(P\); teda o prostyh normalnyh algebrah nad \(P\), ktore kratko budut zvane algebry nad \(P\) ili samo algebry. \(A,\overline A\) sut togda dělitvene algebry. V jeden klas podobnyh algeber \(A,A_f\) sobirajut se vsě \(A_f\) s tym samym prirěđenym \(A\). Izomorfne \(A\) pri tom identificirajut se.

\paragraph{1. Grupa klasov algeber.}
Ako \(A,B\) sut algebry nad \(P\), togda isto važi za njihov produkt, direktny produkt po \(P\); bo po koncu §3 dohodi
\[
        A_r\times_P B_s=C_t,
\]
gde \(C\) jest do izomorfije jednoznačno opreděljena dělitvena algebra s centrom \(P\). Množenje \(A_r,B_s\) s matričnymi kolcami nad \(P\) pokazuje, že to množenje jednoznačno opreděljeno jest za klasy, ktore teda tvoret multiplikativno zakryty komutativny sistem. Kromě togo važi

\paragraph{Teorem R. Brauera.}
Klasy podobnyh algeber tvoret vzhodno k direktnomu tvorjenju produkta abelovu grupu.

Bo sistem ima jediničny klas, sostoječi iz matričnyh kolc nad \(P\), teda algeber podobnyh k \(P\). On ima dalje k vsakomu klasu \((A)\) obratny \((A)^{-1}\), sostoječi iz algeber podobnyh k \(\overline A\), gde \(\overline A\) jest k \(A\) recipročno izomorfno. To slěduje iz komutantnogo teorema §5, 3. posrědstvom specializacije \(S=A\). Bo \(A\) možno nerazložimo vložiti v samo \(A\); dohodi teda \(B=P\) i \(A_A=P_t\).\footnote{Finějše teoremy o grupě klasov algeber upotrěbjajut teoriju faktornyh sistemov; o tom tut ne govorimo. Treba zamětiti, že do togo města ne nastavajut razsmotrenja o komutativnyh poljah; napr. fakt, že rang \((A:P)\) jest kvadrat, nebyl ni upotrěbljen ni dokazan.}

\paragraph{2. Razpadno polje klasa algeber.}
Teorija razpadnyh polj znovu polno opiraje se na princip izrečen na počętku §5; komutativne razširitelne polja predstavjajut se v \(A\) odnosno \(\overline A\). Prědložena bude

\paragraph{Lema.}
Ako \(A\) jest tělo s centrom \(P\), \(Z\) konečno komutativno razširitelno polje od \(P\), togda \(A_Z\) jest prosty s centrom \(Z\). Bo to važi po §4, 1. za k \(A_Z\) recipročno izomorfny sistem \(Z_{\overline A}\). Za vsaky klas \((A)\) podobnyh k \(A\) algeber nad \(P\), \(Z\) porođaje teda klas razširenja \((A)_Z\) podobnyh k \(A_Z\) prostyh algeber nad \(Z\).

Prirěđena dělitvena algebra \(D\) klasa razširenja vyhodi neposrědnje iz komutantnogo teorema §5, 3.:

\paragraph{Teorem o prirěđenoj dělitvenoj algebrě klasa razširenja.}
Nehaj \((A)_Z\) bude klas razširenja, \(Z\) znači nerazložimo vloženje, teda predstavjenje, od \(Z\) v \(A_f\). Togda
\[
        (A)_Z=(D),
\]
gde \(D\) sostoji iz vsih elementov iz \(A_f\), ktore elementno komutujut s \(Z\). Treba samo v §5, 3. zaměniti tamošnji prosty sistem \(S\) črez \(Z\), uzimajuči v obzir, že \(Z_A\) i \(A_Z\) sut recipročno izomorfne. Ako ide o reducibilno vloženje od \(Z\), togda vsě elementy elementno komutujuče s \(Z\) daju matrično kolco nad \(D\).

Komutativno razširitelno polje \(Z\) od \(P\) zove se, kako znano, razpadno polje klasa \((A)\), ako klas razširenja \((A)_Z\) polno razpadaje se, t. j. stavaje ravny jediničnomu klasu nad \(Z\), klasu vsih matričnyh kolc nad \(Z\).

\paragraph{Teorem o harakterizaciji razpadnyh polj.}
Konečno komutativno razširitelno polje \(Z\) od \(P\) jest togda i samo togda razpadno polje klasa \((A)\), kogda jego nerazložimo vloženje v \(A_f\) daje maksimalno komutativno podtělo od \(A_f\).

Bo svojstvo od \(Z\), byti razpadnym poljem od \((A)\), ravnoznačno jest \((D)=(Z)\). Po teoremu o prirěđenoj dělitvenoj algebrě nerazložimo vloženje \(Z\) mora byti maksimalno komutativno podtělo. Ako uslovje izpolnjeno, nužno \(D=Z\), bo vsako \(a\) iz \(D\) črez adjunkciju k \(Z\) porođaje komutativno razširitelno polje.

\paragraph{Zamětky.}
1. Jednočasno slěduje, že maksimalno-komutativno podtělo \(Z\) pri nerazložimom vloženju porođaje maksimalno-komutativno podkolco. Bo vsě elementy elementno komutujuče s \(Z\) tvoret tělo.

2. Teorem daje takodže obstajanje razpadnyh polj; bo prirěđena dělitvena algebra \(A\) sigurano ima maksimalno-komutativne podtěla.

\paragraph{3. Rangove relacije, indeks, beskonečne komutativne razširenja.}
Rangovo tvrdženje iz §5, 4. najprvo daje znany fakt, že rang proste algebry po centru jest kvadrat. Bo ako \(Z\) jest maksimalno-komutativno v \(A\), togda, bo v \(A\) vsako vloženje jest nerazložimo,
\[
        (Z:P)(Z:P)=(A:P),
\]
teda
\[
        (A:P)=m^2,\qquad (Z:P)=m,
\]
gde \(m\) jest Schurov indeks, ktory jednočasno znači stepen vsih v samu dělitvenu algebru \(A\) vložimyh razpadnyh polj. Dalje za stepen \(n\) obćego razpadnogo polja dohodi
\[
        n=mr,
\]
napr. zaradi \((A_r:P)=m^2r^2\), ili takodže po rangovoj relaciji §4, 3., poslednje zato, že indeks \(m\) jest takodže broj prostyh komponentov od \(A_Z\), ktore stavaje matrično kolco nad \(Z\) ranga \(m^2\). Obćejše, ako \(A_L\sim D\), \(L\) jest nerazložimo vloženje od \(A\) v \(A_r\), i \(d^2\) rang \((D:L)\) dělitvenoj algebry \(D\) nad svojim centrom \(L\), togda
\[
        l\,d=mr .
\]
Bo po §5, 4. i teoremu o prirěđenoj dělitvenoj algebrě iz 2. dohodi
\[
        (A_r:P)=(L:P)(D:P),
\]
teda
\[
        m^2r^2=l^2d^2 .
\]

Iz obstajanja razpadnyh polj slědujut za beskonečne komutativne, algebraične ili transcendentne, razširitelne polja \(\Omega\) od \(P\) fakty: klas razširenja \((A)_\Omega\) jest klas prostyh algeber s centrom \(\Omega\). Ako osoblivo \(\Omega\) jest algebraično zakryto, togda \(A_\Omega\) stavaje matrično kolco stepena \(m\). Indeks jest teda jednak ``absolutnomu broju komponentov''. Bo ako \(Z\) jest razpadno polje od \(A\) i \(\Omega'\) kompozitum od \(\Omega\) i \(Z\), togda \(A_{\Omega'}\) jest matrično kolco nad \(\Omega'\), teda bez radikala i s centrom \(\Omega'\). Teda takodže \(A_\Omega\) jest bez radikala i jego centar \(\Omega\), bo od \(\Omega\) različen element od \(A_\Omega\) ležal by takodže v centru od \(A_{\Omega'}\), teda by byl tělo. \(A_\Omega\) jest teda prosta algebra nad \(\Omega\), sigurano matrično kolco nad \(\Omega\), ako \(\Omega\) jest algebraično zakryto.\footnote{Obće obstajajut takodže \glqq{}minimalne\grqq{} razpadne polja, t. j. taka, za ktore žadno pravo podpolje ne jest razpadno polje, vsih stepenovyh čisel. Sr. Brauer--Noether, rabotu citovanu v zam. 2).}

\paragraph{4. Obstajanje separabilnyh razpadnyh polj.}
Vsaky klas \((A)\) ima separabilne razpadne polja, čak taka, ktore možno vložiti v samo \(A\).\footnote{Sr. G. Köthe, \foreign{Über Schiefkörper mit Unterkörpern zweiter Art über dem Zentrum}, Journ. f. Math. 166 (1932), S. 182--184. Tut dany mnogo prostějši dokaz ide nazad k zamětkě M. Zorna.}

Nehaj osnovno polje \(P\) jest karakteristiki \(p\), indeks
\[
        m=s p^f
\]
s \(s\) prěmym k \(p\). Ako \(Z\) jest razpadno polje od \(A\) stepena \(m\) i \(Z_0\) v njem soderženo razširenje prvogo vida, teda
\[
        (Z:Z_0)=p^f,
\]
togda indeks od \(A_{Z_0}\sim D\) ravny jest \(p^f\). Treba dokazati obstajanje separabilnogo razširitelnogo polja \(Z_1\) od \(Z_0\), tako že indeks od \(D_{Z_1}\) stavaje pravi dělitelj od \(p^f\). Ale \(D_\Omega\) s algebraično zakrytym \(\Omega\) po 3. jest bez radikala; reducirana diskriminanta od \(D\) ne izčezava. Teda obstaja najmenej jeden element \(d\) iz \(D\) s neizčezajučim reduciranym slědom; pri tom \(d\) ne može už ležati v osnovnom polju \(Z_0\), bo slěd vsakogo elementa \(\alpha\) iz \(Z_0\) stavaje \(p^f\alpha\), teda izčezava. \(Z_1=Z_0(d)\) stavaje pravo i to separabilno razširitelno polje; indeks od \(D_{Z_1}\) stavaje pravi dělitelj od \(p^f\). Konečno povtorjenje dovodi do separabilnogo razpadnogo polja od \((A)\), čiji stepen \(m\) sovpada s indeksom od \(A\).

\subsection*{§8. Razpadne polja, polja odcěpjenja i Galoisova teorija v komutativnom}

Da by se od razpadnyh polj nad svojim centrom prostyh sistemov prěšlo k razpadnym poljam libovoljnyh prostyh sistemov, treba još razviti razpadnu teoriju centra, k ktorej treba dodati Galoisovu teoriju paralelnu k §6, 3. Vsě těla i sistemy v tom paragrafu sut komutativne.

\paragraph{1. Razpadne i odcěpne polja komutativnyh prostyh sistemov.}
Komutativny sistem \(Z\) nehaj bude prosty nad \(P\), teda polje, \(\Omega\) algebraično zakryto razširitelno polje od \(P\). Togda znano važi: \(Z_\Omega\) ostava sistem bez radikala togda i samo togda, kogda \(Z\) jest separabilno, teda razširenje prvogo vida, nad \(P\). Bo togda i samo togda ima toliko izomorfnyh obrazovanj prvogo stepena v \(\Omega\), koliko jest jego rang po \(P\), teda takodže toliko različnyh predstavjenj, ktore tut sovpadajut s klasami predstavjenj.

Fakt, že v \(Z_\Omega\) može nastati radikal, teda ne nužno ima město direktno sumno razloženje v absolutno proste komponenty, vodi k slědujučim definicijam: razširitelno polje \(\Lambda\) od \(P\) zove se razpadno polje, ako \(Z_\Lambda\) razpadaje se v kompozicione faktory prvogo stepena; drugymi slovami, ako vsě absolutno nerazložime predstavjenja ležet už v \(\Lambda\). Razširitelno polje \(T\) od \(P\) zove se odcěpno polje, ako \(Z_T\) odcěpja najmenej jeden kompoziciony faktor prvogo stepena; drugymi slovami, ako obstaja najmenej jedno absolutno nerazložimo predstavjenje od \(Z\) v \(T\).

Razpadne i odcěpne polja možno harakterizovati po

\paragraph{Teorem.}
Tělo \(T\) jest togda i samo togda odcěpno polje, kogda ono soderži podtělo \(Z'\), izomorfno k \(Z\); tělo \(\Lambda\) jest togda i samo togda razpadno polje, kogda ono soderži podtělo \(\Gamma\), ktore stavaje izomorfno k Galoisovomu tělu naležečemu k \(Z\).

To slěduje direktno iz definicije razpadnyh i odcěpnyh polj posrědstvom absolutno nerazložimyh predstavjenj. Osoblivo vyhodi v analogiji k nekomutativnomu slědstvo: tělo \(Z'\) jest togda i samo togda minimalno odcěpno polje, t. j. nijedno pravo podtělo ne jest odcěpno polje, kogda jest izomorfno k \(Z\). Minimalno odcěpno polje jest togda i samo togda jednočasno razpadno polje, kogda \(Z\) jest normalno, teda Galoisovo tělo nad \(P\).

V tom leži razlog nazvati prostu algebru normalnoju nad svojim centrom. Fakt, že vnutri fiksnogo algebraično zakrytogo \(\Omega\) nad \(P\) obstaja samo jedno jedino minimalno razpadno polje, Galoisovo naležeče k \(Z\), protivno k vobće beskonečno mnogym neizomorfnym v nekomutativnom,\footnote{R. Brauer--E. Noether, a. a. O.} ima osnovu v jednoznačnom direktnom sumnom razloženju v komutativnom proti do operatornoj izomorfije jednoznačnomu v nekomutativnom; ili, čto ravno, v samo konečno mnogyh različnyh absolutno nerazložimyh predstavjenjah vměsto beskonečno mnogyh različnyh v nekomutativnom, ktore vsakože naležet k tomu samomu klasu.

\paragraph{2. Hiperkompleksno obosnovanje Galoisovej teorije v slućaju separabilnyh polj \(Z/P\).}
V točnoj analogiji k §6, 3. tut važi teorija izomorfizmov za libovoljne, nad \(P\) separabilne \(Z\); v slućaju Galoisovogo \(Z\) možno potom prějti k obyčnoj grupě avtomorfizmov.

\paragraph{Predpostavjenja.}
Prosty sistem \(Z\) nad \(P\) nehaj bude konečno separabilno razširitelno polje, \(\Omega\) znači nad \(P\) konečno ili beskonečno razpadno polje, \(Z'=Z^{(1)}\) k \(Z\) izomorfno podtělo, \(\Gamma\) prirěđeno Galoisovo. Po 1. teda važi direktno sumno razloženje
\[
        Z_\Omega=r^{(1)}+\cdots+r^{(n)}
        =e^{(1)}Z_\Omega+\cdots+e^{(n)}Z_\Omega
        =e^{(1)}\Omega+\cdots+e^{(n)}\Omega .
\]
Črez \(e^{(i)}\) porođeno predstavjenje \(Z\to Z^{(i)}\) s \(Z^{(i)}\subseteq\Gamma\) definirano jest črez
\[
        e^{(i)}z=e^{(i)}z^{(i)}.
\]

\paragraph{Lema 1.}
\(n\) predstavjenj porođenyh črez \(e^{(1)},\ldots,e^{(n)}\) sut vsě različne. Dokaz jest kako v §6, 3., ili takodže po 1.; bo one naležet različnym klasam predstavjenj.

\paragraph{Lema 2.}
Ako \(T\) jest medžupolje od \(P\) i \(Z\), \(s\) rang od \(Z\) po \(T\), togda vsaky izomorfizm od \(T\) v \(\Omega\) dopušča najmenej, i zato točno, \(s\) prodolženj.

Dokaz jest kako v §6, 3. Že tut \glqq{}najmenej\grqq{} poostrjaje se k \glqq{}točno\grqq{}, slěduje iz jednoznačnogo razloženja od \(Z_\Omega\), po ktorom \(Z\) ima ne najmenej, ale točno \(n\) izomorfizmov v \(\Omega\).

Kako v §6, 3. definujemo teper črez \(T\) inducirano razděljenje na klasy \(\mathfrak I_T\) konečnoj množiny \(\mathfrak I\) izomorfizmov od \(Z\): vse i samo te izomorfizmy iz \(\mathfrak I\) v \(\mathfrak I_T\) razsmatrajut se kako ekvivalentne, ktore na \(T\) inducirajut toj samy izomorfizm. Dokazuje se

\paragraph{Glavny teorem.}
Ako \(\mathfrak I_T\) jest črez \(T\) inducirano razděljenje na klasy, togda \(T\) jest maksimalno podpolje vzhodno k tomu razděljenju.

Odtud možno za slućaj Galoisovogo \(Z\) točno kako v §6, 3. črez složenje s recipročnym izomorfizmom prějti k grupě avtomorfizmov i k obyčnoj formě. Sootvětno tvrdženje za podgrupy vyhodi iz razsmotrenja idempotentov, ktore samo po sobě jest interesno.

\paragraph{3. Idempotenty od \(Z_\Omega\).}
Nehaj \(S\) proběgaje \(n\) substitucij, ktore prěvodet \(Z\) v pojedina \(Z^{(i)}\), teda složenje izomorfizmov \(Z\to Z_1\) i \(Z\to Z^{(i)}\), pri čem prvi znači recipročny k \(Z\to Z_1\). \(Z\) ne mora byti Galoisovo. Za elementy \(a\) iz \(Z_\Omega\) substitucije \(S\) definirane sut kako operatory črez postavjenje, že na \(Z\) inducirajut identičnost; k vsakomu \(a\) jest teda definirano konjugovano \(a^S\) ležeče v \(Z_{(i)}\). Pri tyh postavjenjah važi

\paragraph{Teorem.}
\(n\) idempotentov \(e^{(i)}\) od \(Z_\Omega\) sut konjugovane:
\[
        e^{(i)}=e^S,\qquad e=e^{(1)};
\]
one ležet v \(n\) konjugovanyh sistemah razširenja \(Z\Omega\).

Bo upotrěbjenje \(S\) prěvodi definicione relacije
\[
        ez=ez^{(1)},\qquad e^2=e
\]
v
\[
        e^S z=e^S z^{(i)},\qquad (e^S)^2=e^S .
\]
Zaradi jednoznačnosti razloženja i idempotenty sut jednoznačno opreděljene; teda \(e^S=e^{(i)}\). Dalje, poneže \(Z\) jest odcěpno polje, teda už posrědkuje predstavjenje \(ez=ez^{(1)}\), \(e\) leži už v \(Z\Omega\), i tym \(e^S\) v \(Z^S\Omega\).

Ako osoblivo \(Z\) jest Galoisovo, slěduje neposrědnje čest glavnogo teorema Galoisovej teorije, odnosna k podgrupam; bo ako \(\mathfrak H\) jest podgrupa Galoisovej grupy \(\Gcal\), zaradi jednoznačnogo opreděljenja komponentov direktnoj sumy element
\[
        \sum_{H\in\mathfrak H} e^H
\]
dopušća vse substitucije iz \(\mathfrak H\), ale nijednu drugu. Poneže grupa \(\Gcal\) definirana byla za \(Z\), ide o Galoisovu teoriju od \(Z\); zaradi izomorfije s tym gotova jest i ta za \(Z'\).

\paragraph{4. Svęz idempotentov s komplementarnymi bazami.}
\(Z\) znovu ne nužno predpostavjaje se kako Galoisovo.

\paragraph{Teorem.}
Ako \(a_1,\ldots,a_n\) jest baza od \(Z\) nad \(P\), i idempotent stavja se
\[
        e=a_1\beta_1+\cdots+a_n\beta_n
\]
teda
\[
        e^S=a_1\beta_1^S+\cdots+a_n\beta_n^S,
\]
togda
\[
        \beta_1^S,\ldots,\beta_n^S
\]
znači obraz komplementarnoj bazy \(b_1,\ldots,b_n\) k \(a_1,\ldots,a_n\) v obrazovanju posrědkovanom črez \(e^S\).

Bo pri obrazovanju črez \(e^S\), \(e^Sz=e^S\zeta^S\), osoblivo \(e^Sa_i=e^S\alpha_i^S\), važe zaradi
\[
        e^Se^S=e^S,\qquad e^Se^R=0\quad(S\ne R)
\]
prirěđenja \(e^S\mapsto1\), \(e^R\mapsto0\). Teda dohodi
\[
        1=\alpha_1^S\beta_1^S+\cdots+\alpha_n^S\beta_n^S,\qquad
        0=\alpha_1^S\beta_1^R+\cdots+\alpha_n^S\beta_n^R\quad(S\ne R).
\]
V matričnom zapisu to jest definiciono uravnenje komplementarnyh baz:
\[
\begin{pmatrix}
\alpha_1 & \cdots & \alpha_n\\
\cdots & \cdots & \cdots\\
\alpha_1^S & \cdots & \alpha_n^S\\
\cdots & \cdots & \cdots
\end{pmatrix}
\cdot
\begin{pmatrix}
\beta_1 & \cdots & \beta_1^R & \cdots\\
\cdots & \cdots & \cdots & \cdots\\
\beta_n & \cdots & \beta_n^R & \cdots
\end{pmatrix}
=
E.
\]
Prěhod k slědovoj definiciji slěduje iz toga, že za
\[
        a_i=\sum_S e^S\alpha_i^S,\qquad b_i=\sum_S e^S\beta_i^S
\]
matrična relacija po zaměně redov prěhodi v
\[
        \Sp(a_i b_i)=1,\qquad \Sp(a_i b_k)=0\quad\text{za }i\ne k,
\]
pri čem \(a_i\) i \(b_i\) ležet v \(Z\), bo dopušćajut vse substitucije \(S\).\footnote{Sr. \foreign{Darstellungstheorie} §25.}\footnote{Svęzi medžu komplementarnymi bazami i idempotentami nahodet se prvi raz pri Dedekindu, \foreign{Zur Theorie der aus \(n\) Haupteinheiten gebildeten komplexen Größen}; sr. Bd. II der Gesammelten Werke, S. 4--5.}

Kontragredientnost komplementarnyh baz slěduje tut iz fakta, že idempotenty \(e^S\) sut invarianty od \(Z\), nezavisne od osnovnoj bazy, teda že vsě
\[
        a_1\beta_1^S+\cdots+a_n\beta_n^S
\]
prěhodet v sebe pri prěhodu k drugoj bazě
\[
        a_1',\ldots,a_n'
\]
od \(Z/P\).

\subsection*{§9. Razpadne polja i odcěpne polja libovoljnyh sistemov}

\paragraph{1. Definicije i redukcija na prosty slućaj.}
Definicijam danym v §8 sootvětsvujut v obćem slućaju definicije: nehaj \(S\) bude hiperkompleksny nad \(P\). Komutativno razširitelno polje \(\Lambda\) od \(P\) zove se razpadno polje od \(S\), ako \(S_\Lambda\) pri kompozicionom rědu iz jednostronnyh idealov razpadaje se v absolutno proste kompozicione faktory; drugymi slovami, ako vsě nerazložime predstavjenja od \(S\) v \(\Lambda\) sut už absolutno nerazložime. Razširitelno polje \(T\) od \(P\) zove se odcěpno polje, ako \(S_T\) odcěpja najmenej jeden absolutno prosty kompoziciony faktor, teda ako najmenej jedno v \(T\) nerazložimo predstavjenje jest už absolutno nerazložimo.

Te definicije jasno soderžet v sebě ty dane v §8 za komutativno i v §7 za proste algebry; poslědnje zato, že tut \(A_\Lambda\) pri vsakoj komutativnoj razširitvi centra jest polno reducibilno i tako kompozicione faktory i komponenty direktnogo sumarnogo razloženja sovpadajut. Polne matrične kolca nad centrom i samo ony daju tut absolutno proste komponenty, teda takodže absolutno nerazložime predstavjenja. Dalje, poneže \(A_\Lambda\) dvostronno prosty jest, teda razpadaje se v operatorno izomorfne komponenty, odcěpjenje absolutno proste komponente už daje polny razpad: odcěpne i razpadne polja sovpadajut.

Iz definicij neposrědnje vyhode fakty: razpadne polja od \(S\) dane sut kako sjedinitelne polja po jednom razpadnom polju v \(P\) nerazložimyh predstavjenj; odcěpno polje od \(S\) jest vsako odcěpno polje jednogo v \(P\) nerazložimogo predstavjenja. Zaradi jedno-jednoznačnogo sootvětsvija prostyh sistemov, komponentov ostatkovogo kolca po radikalu, i v \(P\) nerazložimyh predstavjenj, dosta jest razsmatrjati proste sistemy. Tut rezultaty vyjdut črez kombinaciju §7 i §8.

\paragraph{2. Razpadne polja i odcěpne polja prostyh sistemov.}
Razsmatramo najprvo slućaj prostogo sistema \(A\) se separabilnym centrom \(Z\) nad \(P\). Iz §8, 1. i §8, 3. dohodi: dvostronno razloženje od \(A_\Omega\), sootvětnu direktnomu razloženju od \(Z_\Omega\), gde \(\Omega\) znači razpadno polje od \(Z\), daje \(n\) konjugovanyh, k \(A\) izomorfnyh komponentov
\[
        e^S A_\Omega
\]
od \(A_\Omega\), ktore stavajut proste algebry nad svojim centrom
\[
        e^S Z_\Omega=e^S Z
\]
Pri tom substitucije \(S\) definirane sut kako operatory za \(A_\Omega\) črez postavjenje, že one na \(A\) inducirajut identičnost.

Bo po §8, 3. idempotenty sut konjugovane; isto važi zato za komponenty \(e^S A_\Omega\) po definiciji substitucij za \(A\). Poneže \(A\) dvostronno prosty jest, \(e^S A_\Omega\) kolcovo izomorfny jest k \(A\). Pri tom \(e^SA_\Omega=e^S A_{Z^{S}}\), zaradi \(e^S Z_\Omega=e^S Z\); centar teda sovpada s oblastju koeficientov, i komponenty sut normalne algebry.

Rezultaty iz §7, 2. dalje daju: ako \(A\) jest prosty sistem nad \(P\) s separabilnym centrom \(Z\) nad \(P\), togda takodže \(A_\Omega\) jest bez radikala, teda polno reducibilno; pri tom \(\Omega\) može značiti neko konečno ili beskonečno razširitelno polje. Bo po §7, 2. \(e^S A_\Omega\) ostava prosty pri vsakoj koeficientnoj razširitvi, teda bez radikala. Isto važi za direktnu sumu; ta jest \(A_\Omega\), ili, ako \(\Omega\) ne jest razpadno polje od \(Z\), nastava črez koeficientno razširenje iz \(A_\Omega\).

Dalje iz §7, 2. slěduje

\paragraph{Harakterizacija odcěpnyh polj i razpadnyh polj.}
Ako \(A\) jest dělitvena algebra se separabilnym centrom \(Z\), nerazložime vloženja odcěpnyh polj dane sut vsimi i samo \(Z\) obhvačajučimi maksimalnymi komutativnymi podtělami od \(A_f\); razpadne polja sut vsě i samo sjedinitelne polja odcěpnogo polja s razpadnym poljem centra, osoblivo s Galoisovym tělom naležečim k centru. Minimalno odcěpno polje može byti razpadno polje i bez togo, že centar jest Galoisovy.

Bo odcěpno polje mora obhvačati tělo izomorfno k \(Z\); tvrdženje o vloženju odcěpnyh polj slěduje iz prědhodnyh faktov i iz §7, 2. Razpadno polje mora dalje obhvačati razpadno polje \(\Omega\) za \(Z\). Vsako odcěpno polje nad \(\Omega\) jest pak jednočasno razpadno polje, bo komponenty \(e^S A_\Omega\) sut proste i imajut k \(\Omega\) izomorfny centar. Ako \(Z\) jest Galoisovo, minimalne odcěpne i razpadne polja sovpadajut. Ale i v negaloisovom slućaju može byti, inače než v komutativnom, minimalne odcěpne polja, ktore sut jednočasno razpadne polja, imenno ako tako minimalno odcěpno polje obhvača Galoisovo tělo naležeče k \(Z\). Priměr: kvaternionovo tělo s centrom izomorfnym k \(P(\sqrt[3]{2})\), \(P\) tělo racionalnyh čisel; k \(P(\sqrt[3]{2})\) naležeče Galoisovo tělo jest minimalno odcěpno i razpadno polje.

Ako nakonec ide o neseparabilny centar, togda \(Z_\Omega\) i tym \(A_\Omega\) stavaje sistem s radikalom. Nehaj napr. \(\mathfrak C\) bude radikal od \(Z_\Omega\) i \(\mathfrak c\) ideal razširenja v \(A_\Omega\). Togda za \(Z_\Omega/\mathfrak C\) važi direktno razloženje v konjugovane komponenty, v točnoj analogiji k §8, 2.; to daje kako vyše direktno razloženje od \(A_\Omega/\mathfrak c\), tako že komponenty \(\overline e^{(i)}A\) sut k \(A\) izomorfne s centrom \(\overline e^{(i)}Z^{(i)}\). Teda ostavajut validne i teoremy o odcěpnyh i razpadnyh poljah. Dalje rangovo sčitanje pokazuje, že kompozicione faktory sootvětne kompozicionomu rědu od \(\mathfrak c\) stavajut k \(A_\Omega/\mathfrak c\) operatorno izomorfne, ne samo homomorfne.

Izrečeno za nerazložime predstavjenja, nastavaje rezime: ako dano jest v \(P\) nerazložimo predstavjenje hiperkompleksnogo sistema, togda predstavjenje jest absolutno polno reducibilno togda i samo togda, kogda prirěđeny centar jest separabilny nad \(P\). V vsakom slućaju možno razpadne i odcěpne polja predstavjenja harakterizovati črez vloženje v prirěđeny prosty sistem kako vyše. Osoblivo najnižši stepen odcěpnogo polja ravny jest produktu stepena centra s indeksom prirěđenogo klasa algeber nad centrom; stepen vsakogo odcěpnogo polja jest mnogočlen toga. Predstavjenje razpadaje se, v smyslu kompozicionyh rędov, v toliko različnyh, ale konjugovanyh absolutno nerazložimyh klasov predstavjenj, koliko jest stepen najvelikšego v centru soderžanogo separabilnogo polja. Vsaky klas nastava toliko raz, koliko čini produkt indeksa i stepena centra po tom separabilnom razširenju.

Za sovršeno osnovno polje, gde obstajajut samo separabilne razširitelne polja, tako vrčamo se k znanym rezultatom I. Schura, s dodanjem vložiteljne harakterizacije, ktora se nahodi už pri R. Braueru.

\begin{center}
(Prijeto 8. junija 1932.)
\end{center}

\end{document}
